\documentclass[10pt,a4paper]{amsart}
\usepackage[margin=19mm]{geometry}
\usepackage{amsmath,amssymb,mathtools}
\usepackage[scr=rsfs,cal=euler]{mathalfa}
\usepackage{xfrac}
\usepackage[unicode,hidelinks,bookmarks=false]{hyperref}
\newcommand{\ie}{i.e.\ }
\newcommand{\cf}{cf.\ }
\newcommand{\resp}{respectively\ }
\newcommand{\Subsection}{\subsection}
\providecommand{\nobreakdash}{\nobreak\hbox{-}}
\setlength{\emergencystretch}{2em}
\multlinegap0pt
\usepackage{amssymb}
\usepackage{mathtools}
\usepackage{dsfont}
\usepackage{multirow}
\usepackage{tikz}
\usepackage{booktabs}
\usepackage{bm}
\newtheorem{theorem}{Theorem}
\newcommand{\R}{\mathbb{R}}
\newcommand{\Z}{\mathbb{Z}}
\newcommand{\cA}{\mathcal{A}}
\newcommand{\cM}{\mathcal{M}}
\title{Toric wedge induction and toric lifting property for piecewise linear spheres with a few vertices}
\author{Suyoung Choi, Hyeontae Jang, Mathieu Vall\'ee}
\date{Source: arXiv:2404.15600v1 (CC BY 4.0)}
\begin{document}
\maketitle

\noindent\emph{Source and licence.} Toric wedge induction and toric lifting property for piecewise linear spheres with a few vertices. Version arXiv:2404.15600v1 (CC BY 4.0), distributed by arXiv under the Creative Commons Attribution 4.0 International licence, http://creativecommons.org/licenses/by/4.0/, as recorded in the arXiv licence metadata for this exact version and shown on the abstract page https://arxiv.org/abs/2404.15600v1. Full licence notice: \texttt{licenses/arxiv-2404.15600-CC-BY-4.0.txt}.\par\smallskip

\noindent\textit{Editorial scope.} Counted unit: the article's theorem 168 (source0.tex lines 855--858). The article's complete original proof of it is delivered with it (source0.tex lines 860--1041, one \texttt{proof} environment of 182 lines). No mathematics is written, repaired or re-derived: the statement and the proof are the article's own lines, in the article's own notation, and no retained line is substituted or re-flowed.  No further source-local context is needed: every cross-reference of every retained line resolves inside the two delivered spans.  Nothing was omitted from any retained span, so the package carries no omission record.  No retained line contains a citation, so the wrapper carries no bibliography and no bibliography entry.  No retained line contains a figure environment, an image-inclusion command or drawing code, so no image is delivered and the package declares no asset dependency.  Editorial formatting only, in addition: an AMS \texttt{amsart} preamble that restates the article's own declarations and macros used by the retained lines, namely the environments \texttt{theorem} on separate counters, together with the standard packages those lines need.  The retained environments print in this document as Theorem 1 in order of appearance, whereas the article prints the same items with the numbering of its own full document.  The complete native source of the article as distributed in the arXiv e-print for this exact version is bundled unchanged as \texttt{sources/158082/source0.tex}.
\begin{theorem} \label{theorem: 168}
    Let $K$ be a PL~sphere with less than $168$ facets.
    Then every injective mod~$2$ dual characteristic map over $K$ has a lift as a mod~$2$ characteristic map over $\overline{K}$.
\end{theorem}

\begin{proof}    
    Let $A$ be a $4 \times 15$ matrix over $\Z$ consisting of only $0, 1$ entries with neither repeated columns nor the zero column.
    Consider the binary matroid $\cM$ representing the mod~$2$ linear independence relations between the columns of $A$.
    More explicitly, it is the simplicial complex whose facets are the set of the indices of $4$ columns with an odd determinant in $\Z$.
    Through direct computation, we know that it has $|GL(4,\Z_2)|/4! = 840$ facets, where $GL(4,\Z_2)$ is the general linear group of degree $4$ over $\Z_2$.
    Among them, $835$ correspond to sets of $4$ vectors with determinants $\pm 1$, and the other $5$ correspond to the following sets of $4$ vectors with determinants $\pm 3$:
    \begin{center}
    	$A_1 = \left\{\begin{bmatrix}
    		1 \\
    		1 \\
    		0 \\
    		0
    	\end{bmatrix},
    	\begin{bmatrix}
    		1 \\
    		0 \\
    		1 \\
    		0
    	\end{bmatrix},
    	\begin{bmatrix}
    		1 \\
    		0 \\
    		0 \\
    		1
    	\end{bmatrix},
    	\begin{bmatrix}
    		0 \\
    		1 \\
    		1 \\
    		1
    	\end{bmatrix}\right\}$, 
    	$A_2 = \left\{\begin{bmatrix}
    	1 \\
    	1 \\
    	0 \\
    	0
    	\end{bmatrix},
    	\begin{bmatrix}
    	0 \\
    	1 \\
    	1 \\
    	0
    	\end{bmatrix},
    	\begin{bmatrix}
    	0 \\
    	1 \\
    	0 \\
    	1
    	\end{bmatrix},
    	\begin{bmatrix}
    	1 \\
    	0 \\
    	1 \\
    	1
    	\end{bmatrix}\right\}$, 
    	$A_3 = \left\{\begin{bmatrix}
    		1 \\
    		0 \\
    		1 \\
    		0
    	\end{bmatrix},
    	\begin{bmatrix}
    		0 \\
    		1 \\
    		1 \\
    		0
    	\end{bmatrix},
    	\begin{bmatrix}
    		0 \\
    		0 \\
    		1 \\
    		1
    	\end{bmatrix},
    	\begin{bmatrix}
    		1 \\
    		1 \\
    		0 \\
    		1
    	\end{bmatrix}\right\}$,
    	\medskip
    	
    	$A_4 = \left\{\begin{bmatrix}
    		1 \\
    		0 \\
    		0 \\
    		1
    	\end{bmatrix},
    	\begin{bmatrix}
    		0 \\
    		1 \\
    		0 \\
    		1
    	\end{bmatrix},
    	\begin{bmatrix}
    		0 \\
    		0 \\
    		1 \\
    		1
    	\end{bmatrix},
    	\begin{bmatrix}
    		1 \\
    		1 \\
    		1 \\
    		0
    	\end{bmatrix}\right\}$, and 
    	$A_5 = \left\{\begin{bmatrix}
    		1 \\
    		1 \\
    		1 \\
    		0
    	\end{bmatrix},
    	\begin{bmatrix}
    		1 \\
    		1 \\
    		0 \\
    		1
    	\end{bmatrix},
    	\begin{bmatrix}
    		1 \\
    		0 \\
    		1 \\
    		1
    	\end{bmatrix},
    	\begin{bmatrix}
    		0 \\
    		1 \\
    		1 \\
    		1
    	\end{bmatrix}\right\}$.
    \end{center}
    Let $a_1 = \begin{bmatrix}
    	1 \\
    	1 \\
    	0 \\
    	0
    \end{bmatrix}$, 
    $a_2 = \begin{bmatrix}
    	1 \\
    	0 \\
    	1 \\
    	0
    \end{bmatrix}$, 
    $a_3 = \begin{bmatrix}
    	1 \\
    	0 \\
    	0 \\
    	1
    \end{bmatrix}$, and 
    $a_4 = \begin{bmatrix}
    	0 \\
    	1 \\
    	1 \\
    	1
    \end{bmatrix}$
    be the four vectors in $A_1$.
    In mod~$2$, we can observe that \begin{align*}
    	A_2 &= \{a_1, a_1+a_2, a_1+a_3, a_1+a_4 \}, \\
    	A_3 &= \{a_1 + a_2, a_2, a_2+a_3, a_2+a_4 \}, \\
    	A_4 &= \{a_1 + a_3, a_2 + a_3, a_3, a_3+a_4\}, \text{ and } \\
    	A_5 &= \{a_1 + a_4, a_2 + a_4, a_3+a_4, a_4 \}.
    \end{align*}
    Notice that this combinatorial structure does not depend on the choice of $A_i$'s and $a_j$'s.

    Define $ind_A(A_i) = $ the set of indices of vectors in $A_i$ in $A$.
    For an element $g$ of $GL(4, \Z_2)$, $gA$ is obtained by a column permutation of $A$, so the five sets appear again with some other column indices.
    Suppose that $\{ind_A(A_1), ind_A(A_1), \ldots, ind_A(A_5)\} \cap \{ind_{gA}(A_1), ind_{gA}(A_2), \ldots, ind_{gA}(A_5)\} \not= \varnothing$, that is they contain a common element $ind_A(A_i) = ind_{gA}(A_j)$.
    This means that $g(A_i) = A_j$.
    Then by the property we discussed above, $g$ maps the collection of the five sets $A_i$ on itself.
    Hence for any $g \in GL(4, \Z_2)$, there are only two possibilities:
    \begin{itemize}
    	\item $\{ind_A(A_1), ind_A(A_2), \ldots, ind_A(A_5)\} \cap \{ind_{gA}(A_1), ind_{gA}(A_2), \ldots, ind_{gA}(A_5)\} = \varnothing$ or
    	\item $\{ind_A(A_1), ind_A(A_2), \ldots, ind_A(A_5)\} = \{ind_{gA}(A_1), ind_{gA}(A_2), \ldots, ind_{gA}(A_5)\}$.
    \end{itemize}
    Since every subset of column vectors of $A$ consisting of $4$ vectors with determinant $3$ or $-3$ can be transformed into $A_1$ in mod~$2$ by multiplying with a suitable $g \in GL(4, \Z_2)$, this yields a partition $\cA = \{\{ind_{gA}(A_1), ind_{gA}(A_2), \ldots, ind_{gA}(A_5)\} \mid g \in GL(4, \Z_2) \}$ of the set of facets of $\cM$ with $|\cA| = 840/5 =168$.
    In this partition, only one among the $168$ sets contains the set of vectors of determinant $\pm 3$ in $\Z$.

Now, let $\overline{\lambda^\R}$ be an IDCM over $K$.
    Then, there is an embedding of $\overline{K}$ in $\cM$ according to the index of $\overline{\lambda^\R}(i)$ in $A$ for each vertex $i$.
    If $\overline{K}$ has less than $168$ facets, then there exists an element of $\cA$ which does not intersect the set of the facets of the embedding of $\overline{K}$ by the reverse pigeonhole principle.
    This means that there exists $g \in GL(4, \Z_2)$ such that for any facet $\sigma$ of $\overline{K}$, the determinant of the matrix consisting of the $4$ vectors in $\overline{\lambda^\R} g(\sigma)$ is $1$ or $-1$ when we see the matrix as an integer $\{0, 1\}$-matrix.
    Hence, it provides a lift, as desired.
\end{proof}

\end{document}
